\section{Avance 5: Segunda Alternativa e Integración del Sistema}
\label{sec:avance5}

La segunda alternativa consiste en ejecutar el mismo algoritmo de Dijkstra con el montículo 4-ario indexado (\texttt{DAryHeap}, $d=4$) y con el Fibonacci Heap (\texttt{FibonacciHeap}), las dos estructuras previstas en la Sección \ref{sec:diseno_algoritmico}. Las tres alternativas comparten entradas, formato de salida y criterio de validación. Bellman-Ford actúa como referencia independiente.

\subsection{Interfaz Común e Integración de Componentes}
La función \texttt{solve(instancia, origen, método)} de \texttt{src/algorithms/registry.py} es el único punto de entrada para todos los algoritmos. La interfaz de línea de comandos la expone con los comandos \texttt{list}, \texttt{run}, \texttt{compare}, \texttt{generate}, \texttt{verify}, \texttt{benchmark} y \texttt{menu}.

\begin{table}[H]
    \caption{Módulos reutilizados por todas las alternativas}
    \label{tab:modulos_reutilizados}
    \begin{tabularx}{\textwidth}{@{}>{\raggedright\arraybackslash}p{2.4cm}>{\ttfamily\small\raggedright\arraybackslash}p{4.6cm}X@{}}
        \toprule
        \textbf{Función} & \textbf{Módulo} & \textbf{Descripción} \\
        \midrule
        Lectura & src/io/\allowbreak graph\_io.py & JSON vial (impedancia andina), JSON genérico, lista de aristas y entrada en línea. \\
        Validación & src/io/\allowbreak graph\_io.py & Precondiciones comunes; mensajes con la línea o arista exacta. \\
        Generación & src/generators/\allowbreak instance\_generator.py & Redes aleatorias, cuadrículas y cuadrículas viales con semilla fija. \\
        Ejecución & src/algorithms/\allowbreak registry.py & Selección de alternativa por nombre. \\
        Verificación & src/verification/\allowbreak certificate.py & Certificado de optimalidad y comparación clasificada. \\
        Medición & src/utils/\allowbreak measurement.py & Calentamiento y $n$ pasadas; media, desviación, mínimo y máximo. \\
        Presentación & src/utils/\allowbreak presentation.py & Tablas en consola; exportación y relectura de resultados en JSON. \\
        \bottomrule
    \end{tabularx}
\end{table}

\lstinputlisting[
    language=Python,
    caption={Sift-down del montículo 4-ario: hijo mínimo entre $d$ hijos contiguos (\texttt{python/src/datastructures/d\_ary\_heap.py})},
    label={lst:dary_siftdown},
    firstline=86,
    lastline=104
]{../python/src/datastructures/d_ary_heap.py}

\subsection{Manejo de Entradas Inválidas}
Todas las alternativas comparten la misma validación, de modo que una entrada inválida se rechaza igual sin importar la cola elegida. La CLI devuelve el código 2 ante una entrada inválida y el código 1 si la verificación encuentra un error. Se rechazan: archivos inexistentes, con extensión no soportada o no UTF-8; JSON mal formado; líneas con un número incorrecto de campos; pesos no numéricos, NaN, infinitos o negativos (en Dijkstra); vértices no declarados o duplicados; origen o destino inexistentes; ciclos negativos (en Bellman-Ford); y alternativas desconocidas.

\subsection{Equivalencia de Resultados y Explicación de Diferencias}
Con pesos no negativos, Dijkstra es exacto con cualquier cola de prioridad; por tanto, las distancias de las tres alternativas deben coincidir entre sí y con Bellman-Ford. La función \texttt{compare\_results} clasifica cada diferencia observada (Tabla \ref{tab:clasificacion_diferencias}).

\begin{table}[H]
    \caption{Clasificación de diferencias entre alternativas}
    \label{tab:clasificacion_diferencias}
    \begin{tabularx}{\textwidth}{@{}>{\raggedright\arraybackslash}p{4.6cm}>{\raggedright\arraybackslash}p{2.8cm}X@{}}
        \toprule
        \textbf{Diferencia} & \textbf{Clasificación} & \textbf{Explicación} \\
        \midrule
        Distancia distinta & Error & No debe ocurrir; no se observó en ninguna prueba ni benchmark. \\
        Ruta distinta con igual costo & Propiedad esperada & Existen varios caminos óptimos; $\pi[v]$ solo cambia ante una mejora estricta y cada estructura desempata distinto. Ambas rutas pasan el certificado. \\
        Comparaciones, intercambios, enlaces y cortes & Propiedad de la estructura & Reflejan el diseño de cada cola (Tabla \ref{tab:complejidades}). \\
        \texttt{decrease\_key} y relajaciones (solo con empates) & Propiedad esperada & Dependen del orden de extracción entre vértices de igual distancia (con $|V|=5000$: 6309 frente a 6310). \\
        Dijkstra rechaza y Bellman-Ford resuelve & Criterio distinto & Dijkstra exige $w \ge 0$ (invariante de la Sección \ref{sec:diseno_algoritmico}). \\
        \bottomrule
    \end{tabularx}
\end{table}

\paragraph{Instancia pequeña procesada con todas las alternativas.} Sobre la red del Cusco (métrica distancia), las cuatro alternativas producen distancias y rutas idénticas y pasan el certificado. Las tres colas realizan 6 extracciones, 5 \texttt{decrease\_key} y 9 relajaciones; difieren en las comparaciones (23, 25 y 22) y en sus operaciones propias (5 intercambios en el binario, 3 en el 4-ario, y 3 enlaces más 1 corte en el Fibonacci).

\paragraph{Instancia con empates.} En una cuadrícula $3 \times 3$ de pesos unitarios, las distancias son equivalentes, pero, por ejemplo, para el vértice \texttt{2\_1} (costo 3) Bellman-Ford reporta \texttt{0\_0$\to$0\_1$\to$1\_1$\to$2\_1} y el montículo 4-ario \texttt{0\_0$\to$1\_0$\to$2\_0$\to$2\_1}. Ambas son óptimas.

\paragraph{Instancia propuesta durante la sesión.} El comando \texttt{python main.py run --inline "A B 4; A C 1; C B 2" -s A -q all} (u opción 4 del menú) ejecuta las tres alternativas sobre aristas dictadas en el momento, mostrando la comparación y el certificado.

\subsection{Resultados de Escalabilidad}
\begin{table}[H]
    \caption{Tiempo medio (ms) de Dijkstra por cola sobre redes viales sintéticas (7 pasadas)}
    \label{tab:benchmark_avance5}
    \begin{tabular*}{\textwidth}{@{\extracolsep{\fill}}rrrrrr@{}}
        \toprule
        $|V|$ & $|E|$ & \textbf{Binary} & \textbf{4-ary} & \textbf{Fibonacci} & \textbf{Comparaciones (bin / 4-ary / fib)} \\
        \midrule
        100 & 350 & 0.50 & 0.52 & 0.69 & 1101 / 1347 / 844 \\
        500 & 1750 & 3.17 & 2.98 & 3.85 & 6835 / 8340 / 4853 \\
        1000 & 3500 & 7.42 & 6.84 & 8.94 & 14709 / 17914 / 10197 \\
        2000 & 7000 & 17.49 & 15.39 & 19.52 & 31548 / 37813 / 21202 \\
        5000 & 17500 & 50.94 & 49.11 & 48.80 & 85898 / 100730 / 55948 \\
        \bottomrule
    \end{tabular*}
    \vspace{0.15cm}
    \raggedright \fontsize{9}{11}\selectfont \textit{Nota.} Python 3.13, x86\_64. Grado medio 3.5; semilla $12345+|V|$; 2 pasadas de calentamiento. Antes de medir se verifica la equivalencia de las tres colas.
\end{table}

El Fibonacci Heap realiza entre 23\,\% y 35\,\% menos comparaciones, como predice su análisis amortizado. Sin embargo, es más lento hasta $|V| \approx 2000$ por el costo de crear y enlazar objetos. El montículo 4-ario compara más que el binario, pero hace alrededor de 40\,\% menos intercambios (27\,603 frente a 46\,311 con $|V|=5000$), lo que lo hace igual o más rápido desde $|V|=500$. Con $|V|=5000$ las tres estructuras quedan dentro de una desviación estándar, por lo que no hay un ganador concluyente a esa escala.

\subsection{Pruebas de Integración y Equivalencia}
\texttt{test\_equivalence.py} contrasta las tres colas con Bellman-Ford en 40 grafos aleatorios (pesos enteros con empates, pesos reales y aristas de peso cero), con Dijkstra sobre arreglo y con \texttt{networkx}; además verifica la clasificación de empates y la detección de errores reales. \texttt{test\_cli.py} cubre cada comando, los códigos de salida ante entradas inválidas y el ciclo generar $\to$ ejecutar $\to$ exportar $\to$ verificar. Para comprobar que la suite detecta fallos, se introdujo deliberadamente un error en el \textit{sift-down} del montículo 4-ario: 54 pruebas fallaron.

\subsection{Pendientes, Errores Conocidos y Dificultades}
\textbf{Pendientes:} Pairing Heap, red vial real a mayor escala, medición de memoria y sincronización del cuaderno de Colab. \textbf{Errores conocidos:} no hay errores de corrección abiertos; los tiempos dependen del intérprete y del equipo. \textbf{Dificultades:} separar la equivalencia de distancias (obligatoria) de la de rutas (válida solo si el óptimo es único), y aislar el núcleo del algoritmo de la traza y la reconstrucción de rutas para que no contaminen la medición.
